\section{Uniform expansion before pruning}\label{exp:section}

Fix a height pair $(i,k)$, $|i-k|\leq1$, and fix all the
permutation tables at smaller cutoffs in its family. Throughout this
section write $A=A_i$, $B=B_k$, $D'=D_{ik}$, and
$E=d|A|=D'|B|$. We use the stream conventions, coefficient choices,
and matching coordinates of
Lemma~\ref{tow:coefficient-choice} and
\eqref{tow:edge-equations}; in particular, the independent current
permutations are the cluster matchings described in
Lemma~\ref{tow:matching-clusters}. The graph considered here is $Y^0$:
its current last-block gate is open unconditionally. No assertion of
expansion after pruning is needed.

On each side use uniform probability measure. If $N$ is the slot-average
operator of any realization, then
\[
 (Nf)(a)=\frac1d\sum_{e:\,e_-=a}f(e_+),\qquad
 (N^*g)(b)=\frac1{D'}\sum_{e:\,e_+=b}g(e_-).
\]
Here slots, hence parallel edges, are counted separately. The adjoint
identity follows from $d|A|=D'|B|$. Both operators preserve constants
and are contractions: Jensen's inequality followed by counting slots
gives $\|Nf\|_2^2\leq\|f\|_2^2$, and likewise for $N^*$.
Let $M=\mathbb E N$, with expectation only over the current tables.
The same facts hold for $M,M^*$.

We first bound $M$ on centered functions by organizing them according to
stream prefixes.  The coefficient multiplicity bounds control the full-block
terms; averaging the current tables reduces the terminal contribution to
the partial-block version of the same variance estimate.
We then select actual matchings by a lower bound on every cut and turn
that bound into the required normalized incidence gap.

\paragraph{A conditional variance calculation.}
The following elementary estimate specifies, in particular, the mean
that is subtracted. Its last assertion is needed when the active input
block ends partway through a full block.

\begin{lemma}[Conditional Fourier estimate]\label{exp:fourier}
Let $C$ denote conditioned data. Conditional on $C$, let
$U\in\mathbb F_2^t$ and $V\in\mathbb F_2^m$ be independent uniform
vectors, independent also of auxiliary data $W$. For $1\leq a\leq q$
let $M_a:\mathbb F_2^t\to\mathbb F_2^m$ be linear, let
$o_a=o_a(C,W)$, and let $f_a(C,z)$ be complex functions that do not
depend separately on $U,V,W$. Set
\[
 F=\frac1q\sum_a f_a(C,V+M_aU+o_a),\qquad
 \mu(C)=\frac1q\sum_a2^{-m}\sum_z f_a(C,z).
\]
Suppose that for every nonzero linear functional
$\zeta:\mathbb F_2^m\to\mathbb F_2$, every class of indices on which
$\zeta M_a$ is the same has size at most $\theta q$. Then
\begin{align}
 \mathbb E(F\mid C,W,U)&=\mu(C),\label{exp:conditional-mean}\\
 \mathbb E|F-\mathbb E(F\mid C)|^2
 &\leq\theta\,\mathbb E\left[
    \frac1q\sum_a2^{-m}\sum_z|f_a(C,z)|^2\right].
 \label{exp:conditional-variance}
\end{align}
The conclusion applies with $\theta=10^{-3}$ for a full active
$s$-block and with $\theta=0.55$ for any nonempty active initial
part, using the coefficient multiplicity bounds of
Lemma~\ref{tow:coefficient-choice}.
\end{lemma}

\begin{proof}
Translation of the independent uniform vector $V$ proves
\eqref{exp:conditional-mean}; in particular the right-hand side is
independent of both $U$ and $W$. For fixed $C,W$, write the normalized
Fourier coefficients as
\[
 \widehat f_a(\zeta)=2^{-m}\sum_z
       f_a(C,z)(-1)^{\zeta(z)}.
\]
Orthogonality first in $V$ and then in $U$ gives the exact identity
\[
 \mathbb E_{U,V}|F-\mu(C)|^2
 =\frac1{q^2}\sum_{\zeta\ne0}\sum_{\mathcal C}
 \left|\sum_{a\in\mathcal C}
       (-1)^{\zeta(o_a)}\widehat f_a(\zeta)\right|^2,
\]
where $\mathcal C$ ranges over the classes with a common
$\zeta M_a$. Each square is at most
$|\mathcal C|\sum_{a\in\mathcal C}|\widehat f_a(\zeta)|^2$.
Sum, use Parseval's identity, and integrate over $C,W$. Because the
conditional mean is exactly $\mu(C)$, this integration creates no
additional variance term.
\end{proof}

\paragraph{Prefixes and causal dependence.}
Let $P_h^A$, $h\geq0$, be conditional expectation onto the first
$r+sh$ bits of each of $u_A,v_A$. Let $P_j^B$, $j\geq0$, be
conditional expectation onto the first $sj$ bits of each of $u_B,v_B$.
Only active bits are retained; the others are deterministic zeros.
Thus these are finite increasing filtrations, eventually equal to the
identity, and $P_0^B$ is the projection onto constants. Denote the
constant projection on either side by $C_A,C_B$.

For a fixed $A$-slot and $j\geq1$, the data needed to evaluate a
function in $\operatorname{ran}P_j^B$ are
\begin{equation}
 u_A\big|_{[0,r+s(j-1))},\qquad
 Q_A\big|_{[0,p+sj)}.\label{exp:forward-data}
\end{equation}
Indeed the first $p$ evaluation bits supply $x$, including its bucket,
and inversion of $\mathcal F$ supplies the first $j$ blocks of $u_B$.
At the $j$th $B$-block the keys use only preceding interleaved blocks,
ending with the $(j-1)$st $A$-block. To determine $v_B$ through
length $sj$, the equation for $Q_B$ uses $\mathcal J(u_A)$ through
length at most
\[
 sj-p'=r+s(j-1)-\beta\leq r+s(j-1).
\]
The diagonal and forward-spill rule for $\mathcal K_B$ uses no
$u_B$ beyond its first $j$ blocks to produce these $v_B$ bits.
This proves \eqref{exp:forward-data}, including when a prefix extends
into zero padding.

In the reverse direction, for an input in $\operatorname{ran}P_h^A$
and a fixed $B$-slot, the needed data are
\begin{equation}
 u_B\big|_{[0,sh)},\qquad
 Q_B\big|_{[0,p'+r+sh)}.\label{exp:reverse-data}
\end{equation}
Here $x'$ is read first and the bucket is already part of the slot;
inverting $\mathcal J$ supplies $u_A$ through $r+sh$. Its keys end
with the $h$th $B$-block. Computing $v_A$ through $r+sh$ then uses
$\mathcal F(u_B)$ through at most $sh-\alpha\leq sh$ bits and
the already recovered $u_A$ blocks. These observations prove the
reverse assertion also for $h=0$.

To read a few bits of a diagonal evaluation block one may need its
whole input block. Consequently \eqref{exp:forward-data} and
\eqref{exp:reverse-data} imply, for any fixed tables and also after
taking their mean,
\begin{equation}
 \operatorname{ran}(MP_j^B)\subseteq\operatorname{ran}P_{j+1}^A,
 \qquad
 \operatorname{ran}(M^*P_h^A)\subseteq\operatorname{ran}P_{h+2}^B.
 \label{exp:causality}
\end{equation}
For example $p+sj=r+sj+\alpha$ lies at most $\alpha\leq2$ bits
into the next $A$-block; $p'+r+sh=(h+1)s+\beta$ is handled
by the next $B$-block. This explains both shifts in
\eqref{exp:causality}, including $\alpha=0$ or $\beta=0$, when
the stated bounds can simply be non-sharp.

We now apply Lemma~\ref{exp:fourier} to these data. For
\eqref{exp:forward-data}, condition on $P_{j-1}^A$ and put
$R=r+s(j-1)$. All still needed data from the $A$-state occur through
its evaluation on $[R,R+s+\alpha)$. If its $j$th input block and
this entire $v_A$ interval are active, write that evaluation as
\[
 Z_a=V+M_aU+o_a.
\]
Here $U$ is the full $j$th $A$ input block and $V$ consists of the
$s+\alpha$ new $v_A$ bits. The offset contains the known spill from
the preceding block and any diagonal contribution from the next
input block to the last $\alpha$ output bits. Such next-block input
bits are the auxiliary variable $W$ of Lemma~\ref{exp:fourier}.
They enter nowhere else: \eqref{exp:forward-data} uses direct
$u_A$ only before $R$. Therefore the slot summand is a function
$f_a(C,Z_a)$ with no separate $W$-dependence. The $A$ coefficient
map does not depend on the bucket extracted from $x$; that bucket
may affect $f_a$, but does not change $M_a$. It follows that
\begin{equation}
 \|(1-P_{j-1}^A)MP_j^B\|^2\leq10^{-3}
 \quad\text{whenever these full intervals are active.}
 \label{exp:forward-variance}
\end{equation}

For completeness, the right-hand side of
\eqref{exp:conditional-variance} is bounded by the squared norm of
the original input function, rather than a larger slot-dependent
norm. A uniform output vertex with an independent uniform incident
slot gives a uniform input vertex, since every input has the same
number of slots. Thus the average over slots and vertices of a
deterministic neighbor value squared is exactly the input norm
squared. When a summand averages random table choices, Jensen's
inequality bounds its square by the corresponding average of those
squares. These are precisely the individual norms in
\eqref{exp:conditional-variance}, since $V$ makes each $Z_a$ uniform
conditional on all other data.

The identical argument in \eqref{exp:reverse-data}, conditioning
on $P_h^B$, uses the $B$ input block on $[sh,s(h+1))$ and its
extended evaluation interval $[sh,s(h+1)+\beta)$. It gives
\begin{equation}
 \|(1-P_h^B)M^*P_h^A\|^2\leq10^{-3}
 \quad\text{when that input block and extended $v_B$ interval are full.}
 \label{exp:reverse-variance}
\end{equation}
The explanation of the conditioned mean is the same in both
directions: the means of the slot functions under their uniform
evaluation argument depend only on the retained prefix, so no
uncontrolled contribution comes from the frozen next-block inputs.

\begin{lemma}[Uniform gap for the mean]\label{exp:mean-gap}
For every height pair, and every fixed choice of its lower-cutoff
tables, the mean incidence $M$ satisfies
\[
 \big\|M:\mathbf1_B^\perp\longrightarrow
                   \mathbf1_A^\perp\big\|
 \leq\sqrt{16\cdot10^{-3}+0.55}<0.8.
\]
The bound is uniform in the number of full blocks, the final active
length $1\leq t\leq s$, and $0\leq\alpha,\beta\leq2$.
\end{lemma}

\begin{proof}
Use the orthogonal difference spaces
\[
 E_0^A=P_0^A-C_A,\quad E_h^A=P_h^A-P_{h-1}^A\ (h\geq1),
 \qquad E_j^B=P_j^B-P_{j-1}^B\ (j\geq1).
\]
They decompose the centered spaces; once a filtration saturates,
its further difference spaces are zero. Equation~\eqref{exp:causality}
shows that $E_h^AME_j^B$ can be nonzero only if
\begin{equation}
 h\leq j+1,\qquad j\leq h+2.
 \label{exp:four-bands}
\end{equation}
When $h=j$ or $h=j+1$, the norm of this block is at most
$\sqrt{10^{-3}}$ by \eqref{exp:forward-variance}, provided its
full intervals exist. When $j=h+1$ or $j=h+2$, the same bound
follows by adjointing \eqref{exp:reverse-variance}. For each fixed
value of $h-j$, the input difference spaces are mutually orthogonal,
as are their output difference spaces. Hence a sum along any one
of these four bands has norm at most $\sqrt{10^{-3}}$; their sum
has norm at most $4\sqrt{10^{-3}}$.

There are exactly two forms of terminal cutoff in the interleaved
order. We check the full-interval hypotheses and the remaining
partial-block estimate in each.

\smallskip\noindent
\emph{Case B: the last active block belongs to $B$.}
Write
\[
 l_A=r+ns,\qquad l_B=ns+t,\qquad 1\leq t\leq s.
\]
The active $v$ lengths are
\[
 |v_A|=r+ns+\alpha+t,
 \qquad |v_B|=(n+1)s+\beta.
\]
For $1\leq j\leq n$, the forward estimate uses the full $j$th
$A$ block, and its extended interval ends at
$r+js+\alpha\leq|v_A|$. In the other two bands with $j\leq n$,
one has $h\leq n-1$, so the reverse estimate uses a full
$(h+1)$st $B$ block, ending at $s(h+1)+\beta\leq|v_B|$.
Thus the four-band argument gives
\begin{equation}
 \|M(P_n^B-C_B)\|\leq4\sqrt{10^{-3}}.
 \label{exp:low-B}
\end{equation}

Consider now $M^*g$ for arbitrary $g\in L^2(A)$. Given a $B$-state
and a $B$-slot, inversion through the preceding blocks determines
$u_A$ and all of $v_A$ except its last $t$ bits. In fact all this
determined data use only
\[
 u_B\big|_{[0,ns)},\qquad
 Q_B\big|_{[0,(n+1)s+\beta)}.
\]
The independent current matching averages the remaining $t$
evaluation bits of $Q_A$ uniformly. These bits occupy
$[p+ns,p+ns+t)$, and $\mathcal K_Au_A$ is zero on this interval:
its last possible spill ends at $r+ns+\alpha=p+ns$.
Thus the same bits of $v_A$ are uniformly averaged. Neither their
distribution nor the determined part depends separately on the last
$t$ bits of $u_B$. The tuple defining the matching uses only the
preceding interleaved blocks, the slot, and $x$.

After conditioning on $P_n^B$, the resulting averaged slot function
therefore has the form in Lemma~\ref{exp:fourier}, with $U$ the
$t$ active bits of the last $B$ block and $V$ the full $s+\beta$
bits of $v_B$ starting at $ns$. No current or later input bits
occur except through this evaluation. The multiplicity bound for
the nonempty input prefix, followed by the slot-stationarity and
Jensen argument already given, yields
\begin{equation}
 \|(1-P_n^B)M^*\|^2\leq0.55.
 \label{exp:terminal-B}
\end{equation}
For centered $g$, the two projections $P_n^B$ and $1-P_n^B$ are
orthogonal, and the first component is bounded using the adjoint of
\eqref{exp:low-B}. Consequently
$\|M^*g\|^2\leq(16\cdot10^{-3}+0.55)\|g\|^2$.

\smallskip\noindent
\emph{Case A: the last active block belongs to $A$.}
Write
\[
 l_B=(n+1)s,\qquad l_A=r+ns+t,\qquad1\leq t\leq s.
\]
Now
\[
 |v_A|=r+(n+1)s+\alpha,
 \qquad |v_B|=(n+1)s+\beta+t.
\]
Restrict the $A$ input differences to $0\leq h\leq n$. In the
bands $h=j,j+1$, $j\leq n$, so the forward estimate uses a full
$A$ block and an interval ending by $r+ns+\alpha$. In the bands
$j=h+1,h+2$, $h+1\leq n+1$, so the reverse estimate uses a
full $B$ block and an interval ending by $(n+1)s+\beta$.
All these intervals are active. Thus
\begin{equation}
 \|M^*(P_n^A-C_A)\|\leq4\sqrt{10^{-3}}.
 \label{exp:low-A}
\end{equation}

The current permutation in $\mathcal J$ averages the final $t$
bits of $Q_B$, starting at
$p'+r+ns=(n+1)s+\beta$. The contribution
$\mathcal K_Bu_B$ ends there, so these are exactly the final
$t$ bits of $v_B$. Once they are averaged, a slot value of $Mf$
depends only on
\[
 u_A\big|_{[0,r+ns)},\qquad
 Q_A\big|_{[0,p+(n+1)s)}.
\]
Conditioning on $P_n^A$, use the $t$ active bits of the last $A$
input block and the full $s+\alpha$ evaluation interval beginning
at $r+ns$. Lemma~\ref{exp:fourier} gives
\begin{equation}
 \|(1-P_n^A)M\|^2\leq0.55.
 \label{exp:terminal-A}
\end{equation}
Combining this orthogonally with the adjoint bound from
\eqref{exp:low-A} proves the desired estimate in this case.

The value $t=s$ is included in both arguments. In particular a
cutoff immediately after a full $A$ block is Case A, with $n$ one
less than the number of full $A$ blocks following its initial
$r$ bits; a cutoff immediately after a full $B$ block is Case B.
There is no missing zero-length final block. The actual stream
lengths start with $n_0$ full blocks on each side, so the initial
$r$-bit block is never terminal. If $n=0$ is permitted in the
displayed cases, \eqref{exp:low-B} has zero domain in Case B,
whereas in Case A the only low $A$ difference is $E_0^A$, handled
by \eqref{exp:reverse-variance} with $h=0$. The argument still
applies. When a retained prefix already includes every active bit,
the complementary difference spaces are zero and all the estimates
remain valid.
\end{proof}

\paragraph{Realizing the mean expansion by independent matchings.}
We next choose the current tables so that the actual block has a
uniform edge-isoperimetric constant. A tail estimate for each cut alone
would introduce the number of vertices when summed over all subsets,
and hence would not give the height-uniform bound needed here. Instead,
we will use connected cuts and the local dependence of the matching
variables to select the tables by the local lemma.

Give $A\sqcup B$ the degree measure $\mu(a)=d$, $\mu(b)=D'$;
write $v(Z)=\sum_{z\in Z}\mu(z)$ and
$\mu(A\sqcup B)=2E$. The mean random walk on this Hilbert space is
\[
 R=\begin{pmatrix}0&M\\M^*&0\end{pmatrix}.
\]
The constant-on-each-side subspace has eigenvalues $1,-1$;
its orthogonal complement is the sum of the two sidewise centered
spaces. Lemma~\ref{exp:mean-gap} bounds $R$ above by $0.8$ on
that complement. On the orthogonal complement of the global
constant, the remaining sidewise-constant vector has eigenvalue
$-1$, so the same upper bound $0.8$ holds there. Hence, for
$0<v(Z)\leq E$, the expected number $b(Z)$ of crossing edges obeys
\begin{align}
 \mathbb E b(Z)
 &=\langle(1-R)\mathbf1_Z,\mathbf1_Z\rangle_{\mu}\\
 &\geq0.2\left(v(Z)-\frac{v(Z)^2}{2E}\right)
 \geq0.1v(Z).
 \label{exp:mean-cut}
\end{align}

\begin{lemma}[Tail for a cut]\label{exp:cut-tail}
With lower tables fixed, for every nonempty $Z$ of volume at most
$E$,
\[
 \mathbb P\{b(Z)<v(Z)/20\}\leq\exp(-v(Z)/400).
\]
\end{lemma}

\begin{proof}
Write $X$ for the number of internal matched edges. Then
$b(Z)=v(Z)-2X$. In each independent final cluster, $X$ is a
hypergeometric count: it counts the images in the specified opposite
subset of a specified subset under a uniform bijection. Orient every
cluster from the same side, choosing globally whichever of the
$A$- and $B$-sides of $Z$ has fewer incidences.
Over all clusters the number of sampled incidences is then
\[
 N_Z=\min\{d|Z\cap A|,D'|Z\cap B|\}\leq v(Z)/2.
\]
The cluster counts are independent.

Here is an elementary exponential estimate that also covers clusters
of different sizes. For a sample of $m$ entries without replacement
from a zero-one population with success fraction $\vartheta$, its
exponential moment is at most $(1-\vartheta+\vartheta e^h)^m$.
Indeed this moment is the normalized elementary symmetric polynomial
of order $m$ in the positive entries $1$ and $e^h$.
Replacing any two entries by their average, while keeping all other
entries fixed, increases that polynomial: the only changing term
is their product times a nonnegative elementary symmetric polynomial
of order $m-2$. Repeated averaging and continuity compare with
the list having all entries equal to their mean. For $m=0,1$ the
assertion is immediate. Finally, for a Bernoulli variable $B$,
the second derivative of $\log\mathbb E e^{h(B-\mathbb EB)}$
is a tilted Bernoulli variance, at most $1/4$. Its value and first
derivative vanish at $0$, so
\[
 \mathbb E e^{h(B-\mathbb EB)}\leq e^{h^2/8}
 \qquad(h\in\mathbb R).
\]
Applying these facts cluster by cluster gives
\[
 \mathbb E e^{h(X-\mathbb EX)}\leq e^{N_Zh^2/8}.
\]
By \eqref{exp:mean-cut}, a bad cut requires
$X-\mathbb EX>v(Z)/40$. If $N_Z=0$ the event is impossible.
Otherwise exponential Markov inequality, optimized at
$h=4(v(Z)/40)/N_Z$, bounds its probability by
\[
 \exp\left(-\frac{2(v(Z)/40)^2}{N_Z}\right)
 \leq\exp(-v(Z)/400).
\]
\end{proof}

Let $\mathcal G$ be the finite simple graph of possible endpoints
of edges in this block. Each cluster has at most $2^s$ vertices on
each side. A vertex has at most $\max\{d,D'\}$ incident slots, so
\[
 \Delta(\mathcal G)\leq\Delta:=2D\,2^s,
 \qquad Q:=1+\Delta.
\]
For every nonempty connected vertex set $Z$ of $\mathcal G$ with
$v(Z)\leq E$, introduce the bad event of Lemma~\ref{exp:cut-tail}.
It is enough to avoid these events. Indeed every other set splits
into its connected components in $\mathcal G$; actual edges never
join different components, so the crossing counts and volumes add,
and each component still has volume at most $E$.

Every matching variable is supported on one pair of clusters, with
every edge between their two sides possible. Its support has
$\mathcal G$-diameter at most two. An event for $Z$ depends only
on variables whose supports meet $Z$. Thus two events at distance
greater than two use disjoint sets of independent variables.
Moreover an event is independent jointly of any collection of such
distant events, since all their variables together are disjoint from
its own. We may therefore join two events in a dependency graph
whenever their vertex sets have distance at most two.

There are at most $zQ^2$ vertices within distance two of a $z$-set.
The number of connected $m$-sets containing a specified vertex is
at most $Q^{2(m-1)}$: fix orders of the vertices and their neighbors,
choose for each such set its first spanning tree under that order,
and traverse this rooted tree in depth-first order. The traversal
is a walk of length $2(m-1)$ visiting exactly that set, and the
orders make the association a well-defined injection into such
walks. It follows that a bad event for a $z$-set has at most
\begin{equation}
 zQ^{2m}\quad\text{dependent events indexed by $m$-sets}.
 \label{exp:dependency-count}
\end{equation}
Counting events outside the volume bound or the event itself only
enlarges this upper bound.

Set $a=d/800$ and assign weights $x_Z=e^{-a|Z|}$. For all sufficiently
large $\ell$, uniformly in $i,k$,
\begin{equation}
 a\geq2,\qquad \rho:=Q^2e^{-a}\leq\frac12.
 \label{exp:LLL-parameters}
\end{equation}
Indeed $\log Q=O(\ell)$, whereas $d=2^{4\ell}$; neither bound
contains the stream lengths. Equations~\eqref{exp:dependency-count}
and \eqref{exp:LLL-parameters} give the explicit uniform sum
\begin{equation}
 \sum_{Z'\sim Z}x_{Z'}
 \leq |Z|\sum_{m\geq1}Q^{2m}e^{-am}
 =|Z|\frac{\rho}{1-\rho}\leq|Z|.
 \label{exp:dependency-sum}
\end{equation}
Since every weight is at most $1/2$, the elementary inequality
$\log(1-x)\geq-2x$ yields
\[
 x_Z\prod_{Z'\sim Z}(1-x_{Z'})
 \geq e^{-a|Z|-2|Z|}\geq e^{-2a|Z|}
 \geq\mathbb P(\text{bad }Z).
\]
The last step uses $v(Z)\geq d|Z|$ and
Lemma~\ref{exp:cut-tail}; here $D'\geq D/2\geq d$.

We include the finite asymmetric local-lemma argument for this inequality.
The local lemma originates with Erd{\H o}s and Lov{\'a}sz \cite{ErdosLovasz1975};
see \cite[Theorem~1.1]{MoserTardos2010} for the asymmetric form and
\cite[Section~5.1]{AlonSpencer2016} for its conditional-probability proof. For
any event $A_Z$ and any list $\mathcal S$ of other events, induct
on $|\mathcal S|$ to prove
\[
 \mathbb P\left(A_Z\,\middle|\,
          \bigcap_{A\in\mathcal S}A^c\right)\leq x_Z,
\]
as well as positivity of the avoidance probabilities used. Split
$\mathcal S$ into neighbors $\mathcal N$ of $A_Z$ and nonneighbors
$\mathcal T$. Independence gives
$\mathbb P(A_Z\mid\bigcap_{\mathcal T}A^c)=\mathbb P(A_Z)$.
By the induction hypothesis and the product rule,
\[
 \mathbb P\left(\bigcap_{\mathcal N}A^c\,\middle|\,
                       \bigcap_{\mathcal T}A^c\right)
 \geq\prod_{A_{Z'}\in\mathcal N}(1-x_{Z'}).
\]
If $\mathcal N$ is nonempty, each conditional probability used on
the left conditions on fewer than $|\mathcal S|$ avoided events.
Dividing by this lower bound proves the desired estimate using the
local-lemma inequality above. If $\mathcal N$ is empty, independence
alone proves it. Beginning with the empty list and multiplying
successive survival probabilities also proves their positivity.
In particular the probability of avoiding the entire finite list is
positive. We have obtained an actual block satisfying
\begin{equation}
 b(Z)\geq\frac1{20}v(Z)
 \qquad(0<v(Z)\leq E).
 \label{exp:actual-cut}
\end{equation}

\begin{proposition}[Expansion of the realized incidence]
\label{exp:realized-gap}
For all sufficiently large $\ell$, independently of the height
pair and of its previously fixed lower tables, its current tables
can be chosen so that \eqref{exp:actual-cut} holds. For this choice,
the incidence operator of the actual $Y^0$ block, with each entry
normalized by $\sqrt{dD'}$ in counting Hilbert spaces, sends normalized
constants to normalized constants, as does its adjoint, and has norm
at most $1-\chi$ between their orthogonal complements, where
\[
 \chi=\frac1{800}.
\]
\end{proposition}

\begin{proof}
Only the passage from \eqref{exp:actual-cut} to the incidence norm
remains. We give the weighted discrete Cheeger argument with our
biregular normalization; compare Dodziuk \cite{Dodziuk1984} and
Alon--Milman \cite{AM1985}. Write $R_0$ for the actual random walk, self-adjoint in
degree measure. For real $f$ its Dirichlet form is
\[
 \mathcal E(f)=\langle(1-R_0)f,f\rangle_\mu
             =\sum_{xy\ \mathrm{edge}}(f(x)-f(y))^2.
\]
Each edge in this sum is counted once, with multiplicity. If $w$
is nonnegative and its positive support has volume at most $E$,
integration of \eqref{exp:actual-cut} over the level sets of $w^2$
gives
\[
 \sum_{xy\ \mathrm{edge}}|w(x)^2-w(y)^2|
 =\int_0^\infty b(\{w^2>t\})\,dt
 \geq\frac1{20}\sum_x\mu(x)w(x)^2.
\]
On the other hand Cauchy--Schwarz and
$(w(x)+w(y))^2\leq2(w(x)^2+w(y)^2)$ show that the square of the
left side is at most
\[
 \mathcal E(w)\sum_{xy\ \mathrm{edge}}(w(x)+w(y))^2
 \leq2\mathcal E(w)\|w\|_\mu^2.
\]
Thus $\mathcal E(w)\geq\|w\|_\mu^2/800$, including $w=0$.

For an arbitrary real $f$, choose a degree-weighted median $m$:
both $\{f>m\}$ and $\{f<m\}$ have volume at most $E$.
Such a value is obtained by ordering the finitely many function
values and stopping when cumulative mass reaches half the total.
Apply the preceding estimate to
$w_+=(f-m)_+$ and $w_-=(m-f)_+$. On each edge,
\[
 (w_+(x)-w_+(y))^2+(w_-(x)-w_-(y))^2
 \leq(f(x)-f(y))^2.
\]
This follows by equality when both endpoints are on the same side
of $m$, and by $a^2+b^2\leq(a+b)^2$ otherwise. Consequently,
for a degree-centered $f$,
\[
 \mathcal E(f)\geq\frac1{800}\|f-m\|_\mu^2
 =\frac1{800}\bigl(\|f\|_\mu^2+2E m^2\bigr)
 \geq\frac1{800}\|f\|_\mu^2.
\]
For complex $f$, apply this to its real and imaginary parts. Hence
$R_0\leq1-1/800$ on the orthogonal complement of the global
constant.

Let $T$ be the counting-space incidence divided by $\sqrt{dD'}$.
Multiplication by $\sqrt d$ on $A$ and $\sqrt{D'}$ on $B$ is
a unitary from degree measure to counting measure and sends $R_0$
to
\[
 \begin{pmatrix}0&T\\T^*&0\end{pmatrix}.
\]
Biregularity and $d|A|=D'|B|$ give
$T(|B|^{-1/2}\mathbf1_B)=|A|^{-1/2}\mathbf1_A$ and the
adjoint identity. The centered spaces are invariant in the two
directions. On their direct sum the displayed self-adjoint block
operator has norm $\|T|_{\mathbf1_B^\perp}\|$ and spectrum symmetric
about zero: conjugation by $\operatorname{diag}(1,-1)$ changes its
sign. Its largest spectral value therefore equals this norm.
Every vector in that direct sum is orthogonal to the global
degree constant under the unitary identification, so the preceding
upper gap proves the assertion. This also explains why the global
anti-constant eigenvalue $-1$ does not weaken the incidence bound.
\end{proof}

There are finitely many families. In each family, fix successful
current tables successively in increasing cutoff order. A current
choice does not alter any smaller-cutoff block, whose padded inputs
turn off the higher-cutoff gates. Proposition~\ref{exp:realized-gap}
is uniform over every earlier choice, so this succession defines
all the countably many tables. For example, impose a finite ordering
on each current collection of tables and choose its first successful
element. The construction therefore yields all the required
pre-pruning incidences with the same positive constant $\chi$.
